\chapter{การออกแบบการทดลอง}

\section{ข้อมูลและการสร้าง Label}

ข้อมูลที่ใช้เป็นข้อมูล OHLCV แบบหลายหลักทรัพย์ พร้อม event taxonomy labels เช่น jump, drop, volume spike และ volatility shock การประเมิน dataset ในหลักฐานปัจจุบันชี้ว่าความผิดปกติมี signal เด่นในมิติความผันผวนมากกว่าทิศทางราคาดิบ

\section{การแบ่ง Train, Validation และ Test}

การแบ่งข้อมูลต้องทำก่อน normalization และ scaler ต้อง fit จาก training data เท่านั้น Validation data ใช้สำหรับ threshold, hyperparameter และ model selection ส่วน test data ใช้สำหรับการประเมินสุดท้ายเท่านั้น

\section{Pipeline การทดลอง}

pipeline ปัจจุบันใช้ entry point:

\begin{verbatim}
scripts/gbm/run_joint.py
\end{verbatim}

ซึ่งเรียกกระบวนการ train, validate, test และ visualize ตามลำดับ

\section{Metric}

metric หลักควรรวม precision, recall/sensitivity, specificity, F1, ROC-AUC และ PR-AUC หากมีการเปรียบเทียบหลายรอบทดลอง ควรรายงาน mean, standard deviation, confidence interval และ statistical significance test

